% SIGBOVIK 2026 Submission
% Association for Computational Heresy

\documentclass[sigconf]{acmart}

\usepackage{amsmath}
\usepackage{amssymb}
\usepackage{booktabs}
\usepackage{xcolor}
\usepackage{soul}
\usepackage{microtype}
\usepackage{lipsum}

% For the SIGBOVIK feel
\definecolor{reviewerred}{rgb}{0.8,0.0,0.0}
\newcommand{\reviewer}[1]{\textcolor{reviewerred}{\textbf{[Reviewer \#2: #1]}}}

\title{Negative Vibes Only: Suppressing Unwanted Semantics in Diffusion Models via\\
Aggressive Geometric Algebra and Wishful Thinking}

\subtitle{Also: We Multiplied by 2.0 and It Worked (A Proof)}

\author{Blorptimus Wrengle III}
\affiliation{%
  \institution{Department of Latent Space Janitorial Services}
  \institution{University of Somewhere Cold}
  \city{Somewhere Cold}
  \country{Canada (probably)}
}
\email{b.wrengle@definitely-real-university.edu}

\author{Dr. Perpendicular T. Manifold}
\affiliation{%
  \institution{Institute for Advanced Vibes}
  \city{The Cloud}
  \country{AWS us-east-1}
}
\email{perp@vibes.institute}

\author{Anonymous Reviewer \#2}
\affiliation{%
  \institution{Reviewer \#2's Mom's Basement}
  \city{[REDACTED]}
  \country{[REDACTED]}
}
\email{my-paper-was-rejected@icml.cc}

\begin{abstract}
We present \textit{Normalized Attention Guidance} (NAG), a method for telling
diffusion models what we \emph{don't} want in images, as opposed to the
traditional approach of telling them what we \emph{do} want and just accepting
the disaster that follows. Our key insight is that negative prompts, when applied
in the traditional latent space, are essentially shouting into the void during
the early denoising steps, at which point the model is still a Jackson Pollock
painting and cannot yet be offended by your request to suppress ``ugly, blurry,
watermark.'' Instead, we apply guidance in \emph{attention space}, rotating
feature vectors away from the undesired semantic direction using spherical linear
interpolation (slerp), a technique we borrowed from quaternion animation code
from a 2004 game engine forum post.

Furthermore, we present \textit{BitDance}, a 14-billion parameter model that
generates images 64 tokens at a time in parallel, which we call ``parallel
prediction'' and definitely not ``we ran out of sequential ideas.'' BitDance
uses a text encoder, a projector, a diffusion head, and an autoencoder---which
is either a pipeline or a Rube Goldberg machine depending on whether it works.

We evaluate our methods using FID (which everyone knows is bad but everyone
still uses), a metric we call VIBE (Visually Impressive But Evaluating-nothing),
and a user study conducted on the authors' significant others. Results are
extremely promising assuming you read the right rows of Table 1.

\reviewer{This abstract is too long. Also too short. Reject.}
\end{abstract}

\begin{CCSXML}
<ccs2012>
  <concept>
    <concept_id>10010147.10010257.10010258</concept_id>
    <concept_desc>Computing methodologies~Vibe-based learning</concept_desc>
    <concept_significance>500</concept_significance>
  </concept>
</ccs2012>
\end{CCSXML}

\ccsdesc[500]{Computing methodologies~Vibe-based learning}

\keywords{diffusion models, negative prompts, geometric wishful thinking,
spherical linear interpolation, the number 2.0, latent vibes, attention is all
you need (but guided attention is all you \emph{really} need)}

\maketitle

%-----------------------------------------------------------------------
\section{Introduction}

Text-to-image diffusion models have achieved remarkable quality, enabling users
to generate photorealistic images from natural language descriptions such as
``astronaut riding a horse'' and ``corporate headshot but make it look like I
slept'' \cite{rombach2022ldm,saharia2022imagen,notcited2023something}. A natural
extension is to also specify what one does \emph{not} want: negative prompts
like ``ugly, blurry, deformed, extra fingers'' allow users to describe their
artistic vision as a list of phobias.

Classifier-Free Guidance (CFG) \cite{ho2022cfg} is the standard approach:
\begin{equation}
  \hat{\epsilon} = \epsilon_\theta(x_t, \varnothing) +
    w \bigl(\epsilon_\theta(x_t, c) - \epsilon_\theta(x_t, \varnothing)\bigr)
\end{equation}
which the reader will note involves the empty string $\varnothing$, an
unconditional prediction that is essentially the model's best guess at art
after receiving a prompt equivalent to a blank stare.

\textbf{The problem.} CFG breaks catastrophically in few-step distilled models
like Flux-Schnell, where the model completes the entire creative process in 4
steps (approximately one step per stage of grief). During these steps, especially
the early ones, the conditional and unconditional branches are pointing in
wildly different directions in feature space---imagine asking two blindfolded
strangers to point at a dog. Extrapolating between two random directions does
not produce ``no dogs.'' It produces soup.

\textbf{Our solution.} We observe that attention features are more meaningful
than latent codes early in the denoising process (we observed this by staring
at activation visualizations until they looked like something). We therefore
apply guidance in \emph{attention space} using angular interpolation, ensuring
that our guidance is geometrically meaningful rather than just vibes. The result
is Normalized Attention Guidance (NAG).

\textbf{Also we multiplied by 2.0.} This is correct. We prove this in
Section~\ref{sec:proof_of_two}.

\reviewer{I don't see how this is better than just using more steps. Have you
tried using more steps? I recommend using more steps. Weak reject.}

Our contributions are:
\begin{enumerate}
  \item NAG: guidance in attention space using angular interpolation and
    L1-norm constraints, which reduces artifacts without reducing our
    confidence.
  \item BitDance: a 14B parameter model with a name that will confuse people
    searching for cryptocurrency and TikTok simultaneously.
  \item A formal proof that multiplying by 2.0 at the attention level is
    correct (Section~\ref{sec:proof_of_two}).
  \item Table 1, which contains numbers.
\end{enumerate}

%-----------------------------------------------------------------------
\section{Related Work}

\textbf{Diffusion Models.} There are many papers on diffusion models. We cite
some of them. Others we do not cite, and those authors should probably just
move on.

\textbf{Guidance.} Classifier guidance \cite{dhariwal2021diffusion} trains a
separate classifier to point the denoising process toward desired classes.
Classifier-free guidance \cite{ho2022cfg} eliminates the separate classifier
by training the model to be its own critic, which sounds healthy but isn't.
Perp-Neg \cite{armandpour2023perpneg} applies guidance in the direction
perpendicular to the text embedding, which sounds very sophisticated and
possibly is.

\textbf{Negative Prompts.} Users have discovered that writing ``ugly blurry
deformed'' as a negative prompt generally improves quality, suggesting that
the model has a latent dimension of ``the ugly blurry deformed vibe,'' which
it will add to your image unless explicitly asked not to. This is consistent
with training on the internet.

\textbf{Spherical Linear Interpolation.} Slerp was invented by Ken Shoemake in
1985 \cite{shoemake1985animating} for smooth quaternion animation and has since
been repurposed by every machine learning paper that noticed two vectors and
wanted to interpolate between them. We are no exception.

\textbf{BitDance.} BitDance is related to work on autoregressive image
generation \cite{van2016pixel,esser2021taming}, flow matching
\cite{lipman2022flow}, and that one paper where someone generated images one
row at a time which seemed bad but apparently wasn't. The key difference is
that BitDance does 64 tokens in parallel, which we will now call a
``contribution'' rather than an ``engineering decision.''

\reviewer{You failed to cite [AUTHORNAME 2023], which is my paper. This is
inexcusable. Major revision required.}

%-----------------------------------------------------------------------
\section{Background: Diffusion Models and Why They're Like That}

A diffusion model learns to reverse a noising process. Given a clean image
$x_0$, we add noise:
\begin{equation}
  x_t = \sqrt{\bar\alpha_t}\, x_0 + \sqrt{1-\bar\alpha_t}\, \epsilon,
  \quad \epsilon \sim \mathcal{N}(0, I)
\end{equation}
The model then learns to predict $\epsilon$ (the noise) or $x_0$ (the image)
or $v$ (the velocity), depending on which parameterization the authors
preferred. All three are equivalent, which means three times as many papers
can be written.

\textbf{Flow Matching.} Flow matching \cite{lipman2022flow} is like diffusion
but the noise schedule is straighter, the math is cleaner, and the authors
are more relaxed at conferences. BitDance uses flow matching with velocity
prediction:
\begin{equation}
  v_\theta(x_t, t) \approx x_1 - x_0
\end{equation}
where $x_0$ is noise and $x_1$ is the image, or vice versa, we can never
remember.

\textbf{Transformers.} The transformer architecture \cite{vaswani2017attention}
computes attention weights via:
\begin{equation}
  \text{Attn}(Q, K, V) = \text{softmax}\!\left(\frac{QK^\top}{\sqrt{d}}\right) V
\end{equation}
where $d$ is a number. The output of this attention, before the output
projection $W_o$, is what we call the ``attention hidden state'' $h$,
which is where NAG operates, which is the key thing you need to know.

%-----------------------------------------------------------------------
\section{Normalized Attention Guidance (NAG)}

\subsection{Motivation: The Soup Problem}

As established in Section 1, naive CFG in few-step models produces soup.
Formally, let $h_+$ and $h_-$ be the attention hidden states for positive and
negative prompts respectively. Standard CFG computes:
\begin{equation}
  h_\text{cfg} = h_- + s(h_+ - h_-)
\end{equation}
When $h_+$ and $h_-$ are pointing in random directions (early denoising), this
extrapolation exits the data manifold entirely. The model then decodes this
out-of-manifold vector as soup.

\subsection{Angular Guidance: Rotation Instead of Extrapolation}

Our key insight: instead of extrapolating past $h_+$ in the direction away from
$h_-$, we should \emph{rotate} $h_+$ away from $h_-$. This respects the
manifold structure (citation needed) and produces less soup.

Given $h_+$ and $h_-$, we compute the angular guidance as:
\begin{align}
  \hat{h}_+ &= h_+ / \|h_+\|_2 \\
  \hat{h}_- &= h_- / \|h_-\|_2 \\
  \theta &= \arccos(\hat{h}_+ \cdot \hat{h}_-) \\
  h_\text{ang} &= \frac{\sin(\theta_1)}{\sin(\theta)}\hat{h}_+
                  + \frac{\sin(\theta - \theta_1)}{\sin(\theta)}\hat{h}_-
\end{align}
where $\theta_1$ is chosen to rotate $h_+$ away from $h_-$ by angle
proportional to the guidance scale. The result is then scaled back to the
magnitude of $h_+$:
\begin{equation}
  h_\text{final} = h_\text{ang} \cdot \|h_+\|_2 \cdot 2.0
\end{equation}

The $\times 2.0$ factor is the subject of Section~\ref{sec:proof_of_two}.

\subsection{L1-Norm Guidance: For the Less Geometrically Adventurous}

For those who find spherical geometry unnecessarily exciting, we also provide
an L1-normalized variant:
\begin{align}
  h_\text{cfg} &= h_+ \cdot s - h_- \cdot (s-1) \\
  \text{scale} &= \min\!\left(\frac{\|h_+\|_1}{\|h_\text{cfg}\|_1},\ \tau\right) \\
  h_\text{norm} &= h_\text{cfg} \cdot \text{scale}
\end{align}
where $\tau$ (``tau'') is a hyperparameter the user will need to tune
experimentally, generating approximately 40-60 additional GPU-hours of compute
bills before finding a value that looks good. The final output blends the
normalized guidance with the positive features:
\begin{equation}
  h_\text{out} = \alpha \cdot h_\text{norm} + (1-\alpha) \cdot h_+
\end{equation}
where $\alpha$ is another hyperparameter.

\subsection{The Proof That Multiplying by 2.0 Is Correct}
\label{sec:proof_of_two}

\begin{theorem}[The Factor of Two]
The $\times 2.0$ factor in the angular guidance formula is correct.
\end{theorem}

\begin{proof}
Consider the residual stream of a transformer block. The attention output
$h_\text{attn}$ is added to the residual stream $x$:
\begin{equation}
  x' = x + W_o h_\text{attn}
\end{equation}
NAG modifies $h_\text{attn}$ by a factor of $f$. The effect on the residual
stream is thus scaled by $f$. If $f = 1.25$ (from $\alpha = 0.25$ blending),
and the model has $L = 57$ blocks (Flux has 57 blocks; we counted), then
naive application would result in:
\begin{equation}
  \text{total amplification} = 1.25^{57} \approx 83{,}000
\end{equation}
which is a lot. However, NAG operates \emph{within} the attention computation,
on $h$ before $W_o$, at the attention level rather than the block level.
The residual connection absorbs most of this growth, reducing the effective
amplification to approximately $1.25^1 = 1.25$ per block, which compounds to
$1.25^{57} \approx 83{,}000$\ldots

We pause to reconsider.

Actually the $\times 2.0$ is needed to counteract the $\times 0.5$ implicit in
the attention scale, because the angular guidance reconstructs the vector at
the wrong scale for reasons that are clear from the code if you read it very
carefully and then read it again. Specifically, the slerp output has magnitude
$\|h_+\| / 2$ due to the angular averaging step, and multiplying by $2.0$
restores this. We verified this empirically by running the code with $\times 1.0$
and observing that the images were worse, and then running it with $\times 2.0$
and observing that they were better.

$\square$
\end{proof}

\reviewer{This is not a proof. Weak reject.}

\noindent\textit{Response to Reviewer \#2:} We dispute this characterization.
The square is right there.

%-----------------------------------------------------------------------
\section{BitDance: Dancing with Bits}

BitDance is a 14-billion parameter autoregressive diffusion model designed to
generate images of resolution up to 2048$\times$512, a resolution that is
optimized for people who want very wide images of very short things.

\subsection{Architecture}

The architecture is illustrated in Figure~\ref{fig:arch}, which we did not
actually include in this paper but which you should imagine as a series of
boxes connected by arrows.

\begin{description}
  \item[Text Encoder.] A fine-tuned Qwen3 language model processes the text
    prompt and produces a 5120-dimensional embedding. We use a language model
    as our text encoder because CLIP embeddings felt limiting and we had a Qwen3
    lying around.
  \item[Projector.] A two-layer MLP with GELU activation. This is the part of
    the paper where you normally write ``simple yet effective,'' so: simple yet
    effective.
  \item[Diffusion Head.] A transformer that accepts the projected text features
    and outputs velocity predictions for 64 tokens in parallel. Why 64? Because
    $64 = 8 \times 8$ and $8 \times 8$ is a grid and grids are satisfying.
  \item[Autoencoder.] A convolutional VAE that decodes latent codes into images.
    Uses channel multipliers $(1, 2, 2, 4)$, a sequence so natural that
    stopping at $4$ instead of continuing to $8$ requires genuine restraint.
\end{description}

\subsection{Parallel Prediction: Doing 64 Things at Once}

Standard autoregressive models predict tokens one at a time, autoregressively,
hence the name. This is slow. BitDance predicts 64 tokens simultaneously, which
is fast. We call this ``parallel prediction,'' which is distinct from
``batch processing'' in that it is our idea.

The 64-parallel prediction means BitDance generates a 1024$\times$1024 image
in approximately 15.79 seconds on an A100 GPU, which costs \$2--4 per hour to
rent, meaning each image costs approximately \$0.017. This is approximately
the same as the cost of one ``like'' on social media, a comparison we include
to make the research feel grounded.

\subsection{Special Tokens for Resolution Specification}

BitDance uses special tokens \texttt{<|res\_H|>} and \texttt{<|res\_W|>} to
specify image dimensions. This design decision allows the model to be resolution-
aware, and also means that if you accidentally swap H and W you get a very
wide or very tall image instead of an error, which we have decided to call
a feature.

\subsection{Memory Requirements}

BitDance requires 39--42 GB of GPU memory regardless of resolution, which we
attribute to ``consistent memory allocation patterns'' and definitely not to
the model loading everything into memory immediately and then sitting there.

%-----------------------------------------------------------------------
\section{Experiments}

\subsection{Experimental Setup}

We evaluated our methods on a diverse set of prompts including:
\begin{itemize}
  \item ``A tattered umbrella melts through metallic slats''
  \item ``In a surreal, blurry dreamscape, a massively oversaturated crimson apple''
  \item ``A cyberpunk cityscape at night''
  \item ``Please just make a normal photo'' (negative prompt: ``weird, surreal, dreamscape, oversaturated'')
\end{itemize}

Hardware: one (1) NVIDIA A100 (borrowed). Software: PyTorch, Diffusers, and
approximately 4,000 lines of attention processor code we wrote by hand.

\subsection{Quantitative Results}

\begin{table}[h]
\caption{Comparison of guidance methods. Higher VIBE is better.
Lower FID is better. The $\dagger$ symbol indicates results that are
statistically significant if you don't look too closely at the error bars.}
\label{tab:main}
\begin{tabular}{lccc}
\toprule
Method & FID $\downarrow$ & VIBE $\uparrow$ & Steps \\
\midrule
No guidance & 48.3 & 2.1 & 4 \\
CFG (soup) & 124.7 & 0.4 & 4 \\
NAG (norm\_cfg) & 31.2$^\dagger$ & 8.7$^\dagger$ & 4 \\
NAG (angular) & 29.8$^\dagger$ & 9.2$^\dagger$ & 4 \\
NAG $\times 1.0$ (ablation) & 38.1 & 5.3 & 4 \\
NAG $\times 2.0$ (ours) & 29.8$^\dagger$ & 9.2$^\dagger$ & 4 \\
NAG $\times 3.0$ (curiosity) & 187.4 & 0.1 & 4 & \\
Using more steps (Reviewer \#2) & 32.1 & 7.8 & 28 \\
\bottomrule
\end{tabular}
\end{table}

We observe that NAG outperforms CFG substantially, consistent with our
hypothesis that ``soup'' is bad. The ablation confirms that $\times 2.0$
is correct (see Section~\ref{sec:proof_of_two}). Reviewer \#2's suggestion
(28 steps) does achieve competitive performance but at a 7$\times$ cost
increase, which Reviewer \#2 does not pay.

\subsection{VIBE Metric}

VIBE (Visually Impressive But Evaluating-nothing) is our proposed alternative
to FID. VIBE is computed as follows: the first author looks at the image and
assigns a score from 1 to 10 based on whether they would share it in a group
chat. We believe VIBE captures aspects of image quality that FID misses,
specifically the aspect of ``does it look cool.''

\reviewer{VIBE is not a real metric. Major revision.}

\textit{Response:} FID is also not a real metric by this standard.
We will add a citation to justify VIBE's use.

\subsection{User Study}

We conducted a user study with $n = 3$ participants (the authors). Participants
were shown pairs of images (with and without NAG) and asked to indicate which
they preferred. Results: 100\% preferred NAG images. When the non-first-author
co-author expressed ambivalence about one image pair, their response was recorded
as ``mild preference for NAG'' and the study was concluded.

%-----------------------------------------------------------------------
\section{Discussion}

\subsection{Limitations}

\begin{enumerate}
  \item NAG requires tuning three hyperparameters ($s$, $\tau$, $\alpha$).
    In practice, users will set $\tau = 10$, run the model, think it looks okay,
    and move on. This is fine.
  \item BitDance requires 40 GB of GPU memory, which is the kind of limitation
    that sounds like a limitation but is secretly a ``you should buy an A100''
    advertisement.
  \item The VIBE metric has not been validated against human preferences in any
    systematic way, where ``systematic'' is defined as ``more than 3 people.''
  \item We are not entirely certain the $\times 2.0$ proof is a proof.
\end{enumerate}

\subsection{Broader Impact}

NAG enables more precise negative prompting in text-to-image models. This
technology could be misused to generate images without certain attributes,
such as images without ``watermark'' (legitimate) or images without ``safety
features'' (less legitimate). We therefore recommend that all users behave
themselves, and we acknowledge that this recommendation is insufficient.

BitDance's 64-parallel prediction mechanism could inspire future work in
parallel generation. We claim partial credit for this inspiration.

%-----------------------------------------------------------------------
\section{Conclusion}

We presented NAG and BitDance. NAG operates in attention space, uses angular
interpolation, multiplies by 2.0 (correctly), and produces images without the
features you asked it to avoid. BitDance generates images in parallel batches
of 64 tokens, uses a Qwen3 language model to understand prompts, and costs
approximately \$0.017 per image. Both methods work, as evidenced by Table 1.

Future work includes: (1) explaining why the $\times 2.0$ works in a more
convincing way, (2) reducing BitDance's memory requirements to a quantity that
does not require renting a server, and (3) developing a rigorous theoretical
foundation for the VIBE metric, ideally one that results in a 9.2.

We hope this paper has been informative, or at least entertaining, or at the
very least had enough math in it to look like it might be a real paper at a
glance.

\reviewer{I reviewed a much better paper on a similar topic last year. It was
my own paper. Reject.}

%-----------------------------------------------------------------------
\begin{acks}
The authors thank the SIGBOVIK Program Committee for accepting papers of this
caliber. We also thank Reviewer \#2 for their thorough engagement with the work,
though we note that suggesting ``use more steps'' in response to a paper about
4-step guidance is like suggesting ``use a bigger boat'' in response to a paper
about sailing.

This work was supported in part by GPU hours donated by an institution that
would prefer not to be named in this paper, an A100 that the first author
borrowed ``temporarily'' eight months ago, and several pots of coffee.

No significant others were harmed in the user study, though one expressed
surprise at being asked to rate AI-generated images ``for science.''
\end{acks}

%-----------------------------------------------------------------------
\bibliographystyle{ACM-Reference-Format}
\begin{thebibliography}{99}

\bibitem{rombach2022ldm}
R.~Rombach, A.~Blattmann, D.~Lorenz, P.~Esser, and B.~Ommer.
\newblock High-resolution image synthesis with latent diffusion models.
\newblock In \textit{CVPR}, 2022.

\bibitem{saharia2022imagen}
C.~Saharia et al.
\newblock Photorealistic text-to-image diffusion models with deep language
  understanding.
\newblock In \textit{NeurIPS}, 2022.

\bibitem{notcited2023something}
A.~Somebody.
\newblock A paper we didn't cite but should have.
\newblock \textit{Journal of Things We Forgot}, 2023.

\bibitem{ho2022cfg}
J.~Ho and T.~Salimans.
\newblock Classifier-free diffusion guidance.
\newblock In \textit{NeurIPS Workshops}, 2021.
\newblock (We cited the workshop paper instead of the arXiv version to save
  space in the reference list, not for any other reason.)

\bibitem{dhariwal2021diffusion}
P.~Dhariwal and A.~Nichol.
\newblock Diffusion models beat GANs on image synthesis.
\newblock In \textit{NeurIPS}, 2021.

\bibitem{armandpour2023perpneg}
M.~Armandpour et al.
\newblock Re-imagine the negative prompt algorithm: Transform the negative
  space via geometric principles to enhance coherence and quality.
\newblock \textit{arXiv:2304.05938}, 2023.

\bibitem{shoemake1985animating}
K.~Shoemake.
\newblock Animating rotation with quaternion curves.
\newblock In \textit{SIGGRAPH}, 1985.
\newblock (The only citation in this paper from the 20th century.)

\bibitem{van2016pixel}
A.~van den Oord et al.
\newblock Pixel recurrent neural networks.
\newblock In \textit{ICML}, 2016.

\bibitem{esser2021taming}
P.~Esser, R.~Rombach, and B.~Ommer.
\newblock Taming transformers for high-resolution image synthesis.
\newblock In \textit{CVPR}, 2021.

\bibitem{lipman2022flow}
Y.~Lipman, R.~T.~Q.~Chen, H.~Ben-Hamu, M.~Nickel, and M.~Le.
\newblock Flow matching for generative modeling.
\newblock In \textit{ICLR}, 2023.
\newblock (Published at ICLR 2023 but titled with 2022 because it was on arXiv
  in 2022. This is normal and not confusing.)

\bibitem{vaswani2017attention}
A.~Vaswani et al.
\newblock Attention is all you need.
\newblock In \textit{NeurIPS}, 2017.
\newblock (Every deep learning paper is legally required to cite this.)

\end{thebibliography}

\appendix
\section{Hyperparameter Recommendations}

For users who want NAG to work without reading the paper:

\begin{table}[h]
\caption{Recommended hyperparameters. These were tuned on the authors' preferred
prompts and may or may not generalize.}
\begin{tabular}{lccc}
\toprule
Model & \texttt{nag\_scale} & \texttt{nag\_tau} & \texttt{nag\_alpha} \\
\midrule
Flux-Schnell & 5.0 & 10 & 0.25 \\
Flux-Dev & 5.0 & 2.5 & 0.125 \\
SD3 & 5.0 & 10 & 0.25 \\
SDXL & 5.0 & 10 & 0.25 \\
Wan2.1 (video) & 5.0 & 10 & 0.25 \\
Reviewer \#2's model & N/A & N/A & N/A \\
\bottomrule
\end{tabular}
\end{table}

\section{A Note on the VIBE Metric}

Future work will formalize VIBE as:
\begin{equation}
  \text{VIBE}(x) = \frac{1}{N}\sum_{i=1}^{N} \text{GroupChatShareability}(x, u_i)
\end{equation}
where $u_i$ are human evaluators and \text{GroupChatShareability} is measured
on a Likert scale. We anticipate $N \geq 5$ in the camera-ready version.

\reviewer{Appendix B is not rigorous. Reject.}

\end{document}
